\section*{Hoofdstuk \ref{ch:b2:heart} --- Het hart en de hartcyclus}

\begin{solution}{exo:b2:heart:1}
Holle \omterm{def:b2:blood-circulation:vessels}{ader} $\to$ rechteratrium $\to$ tricuspidalisklep $\to$
rechterventrikel $\to$ pulmonalisklep $\to$ longslagader $\to$ longen
$\to$ longaders $\to$ linkeratrium $\to$ mitralisklep $\to$
linkerventrikel $\to$ aortaklep $\to$ aorta.
\end{solution}

\begin{solution}{exo:b2:heart:2}
Vullen (diastole): AV-kleppen open, halvemaanvormige dicht. Isovolumetrische
contractie: alle vier dicht. Uitstoting: halvemaanvormige open, AV dicht.
Isovolumetrische relaxatie: alle dicht. Eerste toon: de AV-kleppen die aan
het begin van de systole sluiten; tweede: de halvemaanvormige kleppen die aan
het eind ervan sluiten.
\end{solution}

\begin{solution}{exo:b2:heart:3}
Sinusknoop ($t = 0$; de P-golf terwijl de atria depolariseren);
\omterm{prop:b2:heart:conduction}{atrioventriculaire knoop} (bereikt na ongeveer \qty{80}{ms}, houdt de impuls
\qty{100}{ms} vast: het PR-segment); bundel van His en purkinjevezels
(\qtyrange{180}{220}{ms}; het QRS terwijl de \omterm{def:b2:heart:anatomy}{ventrikels} depolariseren); de
T-golf later, wanneer ze repolariseren.
\end{solution}

\begin{solution}{exo:b2:heart:4}
Binnen zijn werkbereik stoot het \omterm{def:b2:heart:anatomy}{ventrikel} meer uit naarmate het meer is
gevuld. Pompt het rechterventrikel meer, dan ontvangt het linker enkele
slagen later meer en pompt het volgens dezelfde wet meer: elke onbalans
corrigeert zichzelf zonder signaal.
\end{solution}

\begin{solution}{exo:b2:heart:5}
\omterm{prop:b2:heart:cycle}{Slagvolume} \qty{80}{mL}; \omterm{prop:b2:heart:cycle}{ejectiefractie} $80/130 = \qty{62}{\%}$;
debiet $80\times 70 = \qty{5.6}{L/min}$. Falend: \qty{60}{mL},
$60/180 = \qty{33}{\%}$, \qty{4.2}{L/min} --- het uitgezette \omterm{def:b2:heart:anatomy}{ventrikel} houdt
het debiet bijna op peil door bij een groot volume te werken, maar stoot maar
een derde uit van wat het bevat.
\end{solution}

\begin{solution}{exo:b2:heart:6}
$W = 100\times 133\times 7\times 10^{-5} = \qty{0.93}{J}$; vermogen
$0.93\times 1.2 = \qty{1.1}{W}$. Bij \qty{150}{mmHg}: \qty{1.4}{J},
\qty{1.7}{W}. De extra \qty{0.56}{W} mechanisch is \qty{3.7}{W}
metabool, \qty{0.19}{mL} zuurstof per seconde: \qty{11}{mL/min} meer.
\end{solution}

\begin{solution}{exo:b2:heart:7}
$60/0.5 = 120$ slagen per minuut; $60/1.5 = 40$. Derde: volledig hartblok
--- de atria kloppen met 100 onder de sinusknoop, de \omterm{def:b2:heart:anatomy}{ventrikels} met 33 onder
een eigen pacemaker, en de AV-knoop geleidt niets.
\end{solution}

\begin{solution}{exo:b2:heart:8}
Stijging $20/25 = \qty{0.8}{s}$, plus \qty{0.15}{s}: \qty{0.95}{s}, 63
slagen per minuut. Acetylcholine: $25/18 = \qty{1.39}{s}$, plus $0.15$:
\qty{1.54}{s}, 39 per minuut. Noradrenaline: $20/40 = \qty{0.5}{s}$,
plus $0.15$: \qty{0.65}{s}, 92 per minuut.
\end{solution}

\begin{solution}{exo:b2:heart:9}
De vertraging laat de atria het leeglopen in de \omterm{def:b2:heart:anatomy}{ventrikels} afmaken voordat
de \omterm{def:b2:heart:anatomy}{ventrikels} samentrekken; zonder haar zouden atria en \omterm{def:b2:heart:anatomy}{ventrikels} samen
samentrekken, zouden de AV-kleppen op half gevulde \omterm{def:b2:heart:anatomy}{ventrikels} dichtgaan en
zou het \omterm{prop:b2:heart:cycle}{slagvolume} dalen met de bijdrage van de atria. Een lang PR-interval
stelt de ventriculaire systole alleen uit en kost niets, tot de vertraging
zo lang wordt dat er slagen uitvallen of de geleiding helemaal faalt.
\end{solution}

\begin{solution}{exo:b2:heart:10}
Rust: $5000/45 = \qty{111}{mL}$; inspanning: $\num{30000}/180 =
\qty{167}{mL}$. De wet van Starling zet de grotere veneuze terugvloed bij
inspanning om in een groter \omterm{prop:b2:heart:cycle}{slagvolume}; de sympathische zenuwen voegen
frequentie en contractiliteit toe (een lager eindsystolisch volume). Boven
ongeveer 200 is de diastole te kort om het \omterm{def:b2:heart:anatomy}{ventrikel} te vullen en daalt het
\omterm{prop:b2:heart:cycle}{slagvolume} sneller dan de frequentie stijgt.
\end{solution}

\begin{solution}{exo:b2:heart:11}
$v = Q/A$: $250/3 = \qty{83}{cm/s}$ en $250/1 = \qty{250}{cm/s}$.
$\Delta P = \tfrac12\times 1060\times v^{2}$: \qty{365}{Pa}
(\qty{2.7}{mmHg}) en \qty{3300}{Pa} (\qty{25}{mmHg}). Het \omterm{def:b2:heart:anatomy}{ventrikel} levert
de extra druk door zijn wand te verdikken; het dikke, stijve \omterm{def:b2:heart:anatomy}{ventrikel}
ontgroeit uiteindelijk zijn kransslagaderlijke aanvoer, vult slecht, en
faalt.
\end{solution}

\begin{solution}{exo:b2:heart:12}
De automatie geeft een slag zonder enige input; de wet van Starling stemt
het debiet af op de terugvloed en houdt de twee \omterm{def:b2:heart:anatomy}{ventrikels} in evenwicht
zonder enige input; de zenuwen en adrenaline stellen alleen frequentie en
contractiliteit bij rond die zelfregelende kern. Een getransplanteerd \omterm{def:b2:heart:anatomy}{hart}
klopt (met ongeveer 100, zonder vagale tonus), stemt zijn \omterm{prop:b2:heart:cycle}{slagvolume} af op
de vulling, en kan zijn debiet bij inspanning verhogen via het
Starlingmechanisme en via circulerende adrenaline --- maar langzaam, en
minder dan een geïnnerveerd \omterm{def:b2:heart:anatomy}{hart}.
\end{solution}

\begin{solution}{pb:b2:heart:1}
\textbf{1.} \qty{70}{mL}; $70/130 = \qty{54}{\%}$; $70\times 70 =
\qty{4.9}{L/min}$.
\textbf{2.} $60/70 = \qty{0.86}{s}$; diastole $0.86 - 0.3 =
\qty{0.56}{s}$.
\textbf{3.} $70/0.3 = \qty{233}{mL/s}$; $233/3 = \qty{78}{cm/s}$.
\textbf{4.} $70/1500 = \qty{0.047}{s}$, ongeveer \qty{50}{ms}.
\textbf{5.} $100\times 133\times 7\times 10^{-5} = \qty{0.93}{J}$;
$0.93\times 70/60 = \qty{1.1}{W}$.
\textbf{6.} Beide \omterm{def:b2:heart:anatomy}{ventrikels} $1.1\times 1.2 = \qty{1.3}{W}$; bij
\qty{15}{\%} \qty{8.7}{W} metabool; $8.7/20 = \qty{0.43}{mL}$
zuurstof per seconde, \qty{26}{mL/min}.
\textbf{7.} Bij een onttrekking van \qty{0.14}{mL} per milliliter \omterm{def:b2:blood-circulation:blood}{bloed}:
$26/0.14 = \qty{190}{mL/min}$.
\textbf{8.} Interval \qty{0.86}{s}, stijgtijd \qty{0.71}{s}:
$20/0.71 = \qty{28}{mV/s}$.
\textbf{9.} Interval \qty{0.6}{s}, stijgtijd \qty{0.45}{s}: \qty{44}{mV/s}.
\textbf{10.} Interval \qty{0.33}{s}, stijgtijd \qty{0.18}{s}: \qty{110}{mV/s}.
\textbf{11.} R-R \qty{0.86}{s}; P: depolarisatie van de atria; QRS:
depolarisatie van de \omterm{def:b2:heart:anatomy}{ventrikels}; T: repolarisatie van de \omterm{def:b2:heart:anatomy}{ventrikels}.
\textbf{12.} De vertraging laat de atria in de \omterm{def:b2:heart:anatomy}{ventrikels} leeglopen voordat
die samentrekken; bij 180 duurt de hele cyclus \qty{0.33}{s}, zodat een
onveranderde vertraging van \qty{0.16}{s} de helft ervan zou opslokken ---
de sympathische zenuwen versnellen ook de geleiding in de knoop.
\textbf{13.} Volledig hartblok: de \omterm{def:b2:heart:anatomy}{ventrikels} met 40 per minuut, in het
ritme gehouden door de bundel of de purkinjevezels, de atria met 75 onder de
sinusknoop.
\textbf{14.} \qty{120}{mL}; $120/150 = \qty{80}{\%}$; $120\times 180 =
\qty{21.6}{L/min}$.
\textbf{15.} $0.33 - 0.3 = \qty{0.03}{s}$: vrijwel geen tijd om te vullen;
nog sneller en het \omterm{prop:b2:heart:cycle}{slagvolume} stort in.
\textbf{16.} $W = 120\times 133\times 1.2\times 10^{-4} = \qty{1.9}{J}$,
$\times 3$ per seconde: \qty{5.7}{W}; beide \omterm{def:b2:heart:anatomy}{ventrikels} \qty{6.9}{W};
metabool \qty{46}{W}; zuurstof \qty{2.3}{mL/s}, \qty{140}{mL/min}.
\textbf{17.} $140/0.14 = \qty{1000}{mL/min}$, vijf keer de rustdoorbloeding,
te leveren in een diastole die tot een tiende van de cyclus is gekrompen: de
\omterm{def:b2:blood-circulation:vessels}{arteriolen} van de kransvaten moeten maximaal verwijden.
\textbf{18.} Frequentie ongeveer 100 zonder vagale tonus, eindsystolisch
volume onveranderd op \qty{60}{mL}: \omterm{prop:b2:heart:cycle}{slagvolume} $150 - 60 = \qty{90}{mL}$,
debiet \qty{9}{L/min}.
\textbf{19.} Frequentie alleen ($70 \to 180$ bij \qty{70}{mL}): $+7.7$;
vulling (\qty{20}{mL} meer bij 180): $+3.6$; contractiliteit
(\qty{30}{mL} minder restvolume bij 180): $+5.4$ --- van 4.9 tot
\qty{21.6}{L/min} (de drie bijdragen, na elkaar bij 180 slagen berekend, tellen
precies op tot het verschil van $16.7$).
\textbf{20.} $233/1 = \qty{233}{cm/s}$, \qty{2.3}{m/s}.
\textbf{21.} $\tfrac12\times 1060\times 2.33^{2} = \qty{2900}{Pa}$,
\qty{22}{mmHg}.
\textbf{22.} $120/0.3 = \qty{400}{mL/s}$, \qty{4}{m/s}; $\tfrac12\times
1060\times 16 = \qty{8500}{Pa}$, \qty{64}{mmHg}.
\textbf{23.} Piek $120 + 22 = \qty{142}{mmHg}$ in rust, ongeveer $140 +
64 = \qty{204}{mmHg}$ bij inspanning; de gemiddelde uitstootdruk in rust
stijgt van 100 tot 122: arbeid per slag $\times 1.22$.
\textbf{24.} Een dikkere wand ontwikkelt meer kracht en verlaagt de spanning
per vezel; maar de te doorbloeden massa groeit, terwijl de
kransslagaderlijke stroom, in de systole dichtgedrukt door de dikke wand en
met een kortere diastole, niet bijblijft, en een dikke wand ontspant slecht
en vult minder: ischemie en falen volgen.
\textbf{25.} \qty{4.9}{L/min} in rust, \qty{21.6}{L/min} bij inspanning;
\qty{0.93}{J} per slag; gradiënt \qty{22}{mmHg} in rust en
\qty{64}{mmHg} bij inspanning.
\end{solution}
